\section{Introduction}

The quest to unite quantum mechanics with general relativity typically culminates in string theory, which posits that our 4D universe is a low-energy effective field theory (EFT) resulting from the compactification of 10D or 11D spacetimes. However, the sheer volume of the string landscape---often estimated at $10^{500}$ possible vacua---has historically precluded deterministic predictions of cosmological parameters. 

Concurrently, empirical observations have revealed persistent tensions in the standard $\Lambda$CDM model, most notably the $S_8$ weak lensing tension observed by KiDS-1000 and DES-Y3, and the Hubble $H_0$ tension. 

In this paper, we propose a radical departure from traditional string landscape statistics. Rather than treating the internal compactification space (modeled here as $K3 \times T^2$) as a random variable, we introduce a dual-track methodology:
\begin{enumerate}
    \item \textbf{Empirical Evolution (AutoEvolve):} A cloud-scale MCMC search optimizing K3 Picard-Fuchs parameters against astrophysical data.
    \item \textbf{Topological Determinism (Wolfram Hypergraphs):} A pure graph-theoretic sieve generating exact algebraic geometries from discrete causal graphs.
\end{enumerate}
Remarkably, both independent tracks converge on identical topological invariants, suggesting a deep, deterministic link between macro-cosmology and micro-topology.
